\documentclass[10pt,a4paper]{article}
\usepackage[T1]{fontenc}
\usepackage{lmodern}
\usepackage[margin=23mm]{geometry}
\usepackage{amsmath,amssymb,booktabs,microtype,tabularx}
\usepackage[hidelinks]{hyperref}
\pdfinfoomitdate=1
\pdftrailerid{}
\pdfsuppressptexinfo=15
\setlength{\parindent}{0pt}
\setlength{\parskip}{5pt}
\begin{document}
\begin{center}
{\Large Tracking the prescribed radiation closure}\\[4pt]
\texttt{exojax-escape}: scoped MCEG review, 23 November 2026
\end{center}

\section{The approximation already in the code}
The current observer model gives every atmospheric layer the same diluted
blackbody mean intensity. It solves local populations with that imposed field:
\begin{equation}
 J_\nu(r)=W B_\nu(T_{\rm rad}),\qquad
 \bar n_\nu(r)=\frac{c^2J_\nu(r)}{2h\nu^3}
 =\frac{W}{\exp[h\nu/(k_{\rm B}T_{\rm rad})]-1}.
 \tag{RAD-1}\label{eq:radiation-closure}
\end{equation}
Here $W$ is dilution, $T_{\rm rad}$ is radiation temperature, and $\bar n_\nu$
is photon occupation. The field is independent of layer radius $r$. This is
Equation \texttt{radiation\_closure@1} in
\href{../mceg/contracts/escape.nlte_spectrum.yaml}{the spectrum contract}.
Binding \texttt{prescribed\_radiation} identifies the
\texttt{population\_profiles} method of \texttt{HydrogenNLTEForwardModel}
in \texttt{src/escape/nlte\_forward.py}. It computes line and continuum
occupations once and shares them across layers.

Observer transfer computes emergent intensity but does not return attenuated
or internally generated radiation to the atom. The approximation was documented
when this implementation was introduced; its quantitative error was unknown.

\section{Why it needs a separate scientific question}
Zhang et al. use a PHOENIX incident spectrum and Cloudy to obtain temperature
and ionization structure (\href{https://arxiv.org/html/2606.23899v1\#S4.SS1}{Section IV.1}).
Cloudy distinguishes the incident boundary field, its attenuation, and diffuse
radiation (\href{https://www.pa.uky.edu/~gary/CloudyWorkshop/CloudyWorkshop/16Shandong/Cloudy_files/hazy1.pdf}{Hazy 1, Section 2.2}).
The complete input configuration of the paper's c23.01 run has not been recovered.

A later column diagnostic gives Lyman-edge optical depths of about 19 and 41
at two representative H-alpha layers; static Lyman-alpha depths are about
$1.8\times10^5$ and $3.2\times10^5$. These \emph{prescribed-field} columns
rule out an optically thin justification there, but do not determine the actual
mean intensity or H-alpha error. The
\href{../mceg/evidence/radiation_review_20261123/review.json}{compact review}
retains exact values and immutable archive/member hashes. This update runs
no new physical calculation.

\section{What MCEG now keeps connected}
\begin{tabularx}{\linewidth}{@{}lX@{}}
\toprule
Record & Meaning for this approximation \\
\midrule
Claim & \texttt{prescribed\_field\_adequacy}: is its error acceptable for the intended use? \\
Assessment & \texttt{prescribed\_field\_review}: \textbf{unknown}. The required comparison and physical error budget are missing. \\
FollowUp & \texttt{test\_radiation\_feedback}: perform and review the controlled comparison below. \\
Decision & \texttt{prescribed\_field\_use}: retain exploratory use; physical inference is not endorsed. \\
\bottomrule
\end{tabularx}

Hold density, temperature, incident spectrum, atomic data and observer settings
fixed. Compare the imposed field with radiation transport coupled to statistical
equilibrium, including a matched thin limit and the retained thick-column cases.
Record mean intensities, ionization, $n_2,n_3$ and the same H-alpha observable.
Check numerical convergence separately. Specify the physical error budget and
parameter domain before interpreting the difference; the existing $0.1\sigma$
numerical refinement criterion is not automatically a model-error allowance.

Completion requires retained evidence and scoped Assessments, with a next
action or restricted-use Decision for remaining issues. The controlled numerical
comparison gets its own numerical Claim; it does not alone establish physical
adequacy. Solver convergence or a Decision does not turn \texttt{unknown} into
\texttt{supported}. FollowUp \texttt{validate\_observer\_application} separately
tracks observation processing and the other atmospheric limitations.

The earlier \href{nlte_spectrum.pdf}{observer PDF} and its evidence retain their
hashes. New Equation/Binding pins identify this approximation without validating
it. The decorated JAX method still yields an unresolved static declaration
check. MCEG remains external to the runtime package \texttt{escape}.
\end{document}
